%! Author = chtompki
%! Date = 1/14/26
% Example LaTeX document for GP111 - note % sign indicates a comment
\documentclass[12pt]{amsart}
% Default margins are too wide all the way around. I reset them here
\setlength{\hoffset}{-.75in}
\setlength{\textwidth}{6.5in}
\setlength{\voffset}{-.5in}
\setlength{\textheight}{9.0in}
\usepackage{amsmath}
\usepackage{amssymb}
\usepackage{amsthm}
\usepackage{pgfplots}
\usepackage{enumerate}
\usepackage{tikz}
\usetikzlibrary{arrows.meta} % Required for the Circle and Triangle arrow tips
\usepackage{setspace}

\theoremstyle{definition}
\newtheorem{definition}{Definition}

\title{Seminar Summary:\\February 10, 2026 - David Chan\\Human/Individual Research}
\author{Rob Tompkins}
\date{\today}


\begin{document}
    \maketitle
    \date{\today}
    \begin{onehalfspace}
        The goal behind the talk by Professor Chan on February $10^{\text{th}}$ was to explore three new research
        topics that we did not go over in his initial talk on January $12^{\text{th}}$.
        More specifically we covered the research topics: \emph{Contact Tracing}, \emph{Student Support Networks}, and
        \emph{Instructor Student Interaction}. We will explore the different research projects below in the subsequent
        sections.

        \section{Contact Tracing}
        Contact-tracing is the physical process of retracing human interconnectivity over a brief window of time in hopes
        of determining a full graph of distinct people that a single individual had contact with during that window of time
        to help determine and mitigate potential disease spread.
        During the Covid pandemic we saw remedial efforts at contact tracing in the population, but because the
        disease's symptoms were so mild during the period of contagion for an individual and the symptoms of the
        disease were so mild for a considerable portion of the population we were largely unable to do anything about
        it via these processes.
        Where we see hugely effective contact tracing is in populations where the disease is far more leathal and the
        mechanism of transfer between patients is non-airbourne and more based in physical contact.
        For example in Africa, the infrastructure and knowledge around contact tracing due to Ebola is profoundly better
        than anywhere else in the world.



    \end{onehalfspace}
\end{document}
